\begin{table}[htbp]
\centering\small
\caption{M2 large-firm contrasts under alternative clustering}
\label{tab:status-clustering}
\begin{tabular}{l ll rrr}
\toprule
Status & Reference & Formality & Premium (\%) & $p$ (SD) & $p$ (S) \\
\midrule
Private-sector employees & Solo & Formal & 16.2 & $<0.001$ & $<0.001$ \\
Private-sector employees & Solo & Informal & 16.4 & $<0.001$ & 0.014 \\
Own-account workers & Solo & Formal & 53.9 & $<0.001$ & 0.001 \\
Own-account workers & Solo & Informal & 46.2 & $<0.001$ & $<0.001$ \\
Private-sector employees & 11-19 & Formal & 13.1 & $<0.001$ & $<0.001$ \\
Private-sector employees & 11-19 & Informal & -4.3 & 0.002 & 0.055 \\
Own-account workers & 11-19 & Formal & 24.4 & $<0.001$ & 0.004 \\
Own-account workers & 11-19 & Informal & 12.1 & $<0.001$ & 0.073 \\
Government employees & 11-19 & Formal & 17.6 & $<0.001$ & 0.012 \\
Government employees & 11-19 & Informal & -23.0 & 0.110 & 0.348 \\
Employers & 11-19 & Formal & 57.5 & $<0.001$ & 0.019 \\
Employers & 11-19 & Informal & 36.2 & 0.077 & 0.054 \\
Day laborers & Solo & Formal & 3.1 & 0.441 & $<0.001$ \\
Day laborers & Solo & Informal & 20.1 & $<0.001$ & $<0.001$ \\
\bottomrule
\end{tabular}
\par\vspace{0.3em}
\begin{minipage}{0.97\textwidth}\footnotesize
Notes: SD denotes sector-department clustering and S denotes sector clustering. Coefficients and samples do not change. Day laborers are concentrated in agriculture. Their reference cells span very few sectors, limiting the informativeness of sector-clustered inference.
\end{minipage}
\end{table}
